\documentclass[11pt,a4paper]{article}

\usepackage[brazil]{babel}
\usepackage[utf8]{inputenc}
\usepackage[T1]{fontenc}
\usepackage[a4paper,margin=1.8cm]{geometry}
\usepackage{tabularx}
\usepackage{longtable}
\usepackage{array}
\usepackage{amssymb}
\usepackage[sfdefault]{atkinson}
\usepackage{multicol}
\usepackage{fancyhdr}
\usepackage{lastpage}
\usepackage{titlesec}
\usepackage{booktabs}
\usepackage{xcolor}
\usepackage{framed}
\usepackage{needspace}

% ---------- Paleta ----------
% Um preto levemente suavizado para títulos (menos duro que o preto puro) e
% dois cinzas para linhas/campos, para reduzir o peso visual das bordas sem
% comprometer a legibilidade em impressão P&B.
\definecolor{ink}{RGB}{38,38,38}
\definecolor{hairline}{RGB}{140,140,140}
\definecolor{fieldline}{RGB}{110,110,110}

\titleformat{\section}  % which section command to format
  {\color{ink}\fontsize{12}{12}\bfseries} % format for whole line
  {\thesection} % how to show number
  {1em} % space between number and text
  {} % formatting for just the text
  [{\color{hairline}\titlerule[0.5pt]}] % thin rule under the whole heading
\titlespacing*{\section}{0pt}{1.4em}{0.6em}

\setlength{\parindent}{0pt}
\renewcommand{\arraystretch}{1.4}

% ---------- Campos de preenchimento ----------

\newcommand{\digitbox}{
    {\color{fieldline}\fbox{\color{black}\rule{0pt}{2.2ex}\hspace{0.7em}}}
}

\newcommand{\twodigits}{
    {\color{fieldline}\fbox{\color{black}\rule{0pt}{2.2ex}\hspace{0.7em}}}\null
    {\color{fieldline}\fbox{\color{black}\rule{0pt}{2.2ex}\hspace{0.7em}}}
}

\newcommand{\threedigits}{
    {\color{fieldline}\fbox{\color{black}\rule{0pt}{2.2ex}\hspace{0.7em}}}\null
    {\color{fieldline}\fbox{\color{black}\rule{0pt}{2.2ex}\hspace{0.7em}}}\null
    {\color{fieldline}\fbox{\color{black}\rule{0pt}{2.2ex}\hspace{0.7em}}}
}

\newcommand{\timefield}{
    \twodigits {\Large\textbf{:}} \twodigits
}

\newcommand{\datefield}{
    \twodigits {\Large\textbf{/}} \twodigits {\Large\textbf{/}} \threedigits
}

\newcommand{\dotted}[1]{
    \color{fieldline}\makebox[#1]{\dotfill}\color{black}
}

% Linha numerada para anotação livre (usada em "Ocorrências").
\newcommand{\notedline}[1]{
    \noindent{\color{ink}\bfseries #1.}\ {\color{fieldline}\rule[-0.1cm]{0.9\textwidth}{0.4pt}}\par
    \vspace{0.9cm}
}

% Chamada lateral (usada nas instruções de "Registro de Presença") — uma
% barra fina à esquerda em vez de uma caixa fechada, para um visual mais leve.
\renewcommand{\FrameCommand}{ {\color{fieldline}\vrule width 1.6pt}\hspace{10pt}}
\setlength{\OuterFrameSep}{0pt}
\setlength{\FrameSep}{4pt}

\pagestyle{fancy}
\fancyhf{}
\fancyhead[L]{\color{ink}\textbf{ATA DE APLICAÇÃO DE PROVA}}
\fancyhead[R]{\color{ink}UFRGS}
\fancyfoot[L]{\small\color{ink} Gerado em {{ generated_at }}}
\fancyfoot[R]{\color{ink}Folha \thepage \ de \pageref{LastPage}}
\renewcommand{\headrulewidth}{0.5pt}
\renewcommand{\footrulewidth}{0.5pt}
\renewcommand{\headrule}{ {\color{hairline}\hrule width\headwidth height\headrulewidth}}
\renewcommand{\footrule}{ {\color{hairline}\hrule width\headwidth height\footrulewidth}}

\begin{document}

\begin{center}
    {\Large\bfseries\color{ink} Aplicação de Prova}
\end{center}

% ====================================================
% IDENTIFICAÇÃO
% ====================================================

\begin{tabularx}{\textwidth}{lXlX}
\textbf{Atividade de Ensino}
    & \multicolumn{3}{l}{ {{ exam.course }} } \\

{% if exam.features.course_code.is_some() %}
\textbf{Código}
    & \multicolumn{3}{l}{ {{ exam.features.course_code.as_ref().unwrap() }} } \\
{% endif %}

\textbf{Período Letivo}
    & \multicolumn{3}{l}{ {{ exam.semester }} } \\

\textbf{Professor Responsável}
    & \multicolumn{3}{l}{ {{ exam.professor }} } \\

{% if exam.features.department.is_some() %}
\textbf{Instituto/Departamento}
    & \multicolumn{3}{l}{ {{ exam.features.department.as_ref().unwrap() }} } \\
{% endif %}

\textbf{Turma(s)}
    & \multicolumn{3}{l}{ {{ exam.classes }} } \\

\textbf{Avaliação}
    & {{ exam.exam_name }}
&
\textbf{Aplicação}
    & {{ exam.application_type }} \\

\textbf{Data}
    & \multicolumn{3}{l}{ {{ exam.date_br() }} } \\

\textbf{Horário Previsto}
    & \multicolumn{3}{l}{
        {{ exam.schedule.scheduled_start_br() }}
        --- {{ exam.schedule.scheduled_end_br() }}
        ({{ exam.schedule.duration_minutes }} minutos)
    } \\

\textbf{Sala}
    & \multicolumn{3}{l}{ {{ exam.room }} } \\

{% if exam.features.allowed_materials.is_some() %}
\textbf{Materiais Permitidos}
    & \multicolumn{3}{l}{ {{ exam.features.allowed_materials.as_ref().unwrap() }} } \\
{% endif %}
\end{tabularx}

% ====================================================
% PRESENÇA
% ====================================================

\section*{Registro de Presença}

\begin{leftbar}
\small
Marque somente os alunos presentes com \textbf{P} e os ausentes com \textbf{A} na
coluna de presença no início da aplicação. Os alunos que chegarem após o início da
aplicação devem ser marcados com \textbf{T} (tarde).
\par\smallskip
Os alunos devem assinar na coluna de assinatura no ato da devolução da prova.
\end{leftbar}

\medskip

\begin{longtable}{ {{ exam.features.presence_table_colspec()|safe }} }

\toprule
\textbf{Nº} & \textbf{Matrícula} & \textbf{Nome}{% if exam.features.show_class_column %} & \textbf{Turma}{% endif %} & \textbf{Presença} & \textbf{Assinatura}{% if exam.features.show_delivery_time_column %} & \textbf{Entrega}{% endif %} \\
\midrule
\endfirsthead

\multicolumn{ {{ exam.features.columns_after_number() + 1 }} }{l}{\small\textit{(continuação da lista de presença)}} \\[0.3em]
\toprule
\textbf{Nº} & \textbf{Matrícula} & \textbf{Nome}{% if exam.features.show_class_column %} & \textbf{Turma}{% endif %} & \textbf{Presença} & \textbf{Assinatura}{% if exam.features.show_delivery_time_column %} & \textbf{Entrega}{% endif %} \\
\midrule
\endhead

\midrule
\multicolumn{ {{ exam.features.columns_after_number() + 1 }} }{r}{\small\textit{continua na próxima página\ldots}} \\
\endfoot

\bottomrule
\endlastfoot

{% for student in exam.students %}
{{ student.number }}
&
{{ student.id }}
&
{{ student.name }}
{%- if student.observations.is_some() %}
\newline{\footnotesize\textit{Obs.: {{ student.observations.as_ref().unwrap() }}}}
{%- endif %}
{% if exam.features.show_class_column %}
&
{{ student.class }}
{% endif %}
&
$\digitbox$
&
\dotted{\linewidth}
{% if exam.features.show_delivery_time_column %}
&
\dotted{\linewidth}
{% endif %}
\\[0.3em]
{% endfor %}

\end{longtable}

% ====================================================
% HORÁRIOS
% ====================================================

\section*{Horários}

\begin{tabularx}{\textwidth}{lrlr}

Horário previsto de início:
&
{\Large\textbf{ {{ exam.schedule.scheduled_start_br() }} }}
&
Horário real de início:
&
\timefield
\\

Tempo de prova:
&
{\Large\textbf{ {{ exam.schedule.duration_minutes }} minutos }}
&
&
\\

Horário limite previsto:
&
{\Large\textbf{ {{ exam.schedule.scheduled_end_br() }} }}
&
Horário limite calculado:
&
\timefield
\\

&
&
Horário real de término:
&
\timefield

\end{tabularx}

% ====================================================
% SAÍDAS TEMPORÁRIAS
% ====================================================

\section*{Registro de Saídas Temporárias}

\begin{tabularx}{\textwidth}{XccX}

\toprule
Nº Aluno &
Saída &
Retorno &
Rubrica do Fiscal \\
\midrule

{% for _ in 0..exam.exit_rows %}
\dotted{4.5cm}
&
\timefield
&
\timefield
&
\dotted{5.5cm}
\\
{% endfor %}

\bottomrule
\end{tabularx}

% ====================================================
% OCORRÊNCIAS
% ====================================================

\section*{Ocorrências}

\small\textit{Descreva abaixo qualquer ocorrência relevante durante a aplicação (fraude,
mal-estar, atraso, saída antecipada etc.), indicando horário e aluno(s) envolvido(s),
quando aplicável.}

\bigskip

\begin{minipage}{\textwidth}
\notedline{1}
\notedline{2}
\notedline{3}
\notedline{4}
\notedline{5}
\end{minipage}

% ====================================================
% ENCERRAMENTO
% ====================================================

% O bloco de encerramento tem altura fixa (não cresce com o número de
% alunos), então evitamos que ele seja partido ao meio por uma quebra de
% página — se não couber inteiro, ele inteiro vai para a próxima página.
\needspace{9cm}

\section*{Encerramento}

Total de alunos matriculados na lista de presença: \textbf{ {{ exam.students.len() }} }

\begin{multicols}{2}
    \begin{tabularx}{\textwidth}{lX}
        Total de presentes:         &
        \twodigits
        \\[0.2cm]

        Total de ausentes:          &
        \twodigits
        \\[0.2cm]

        Provas recolhidas:          &
        \threedigits
        \\[0.2cm]

        Folhas extras distribuídas: &
        \twodigits
        \\[0.2cm]

        Folhas extras recolhidas:   &
        \twodigits
    \end{tabularx}

    \textbf{
        Conferência final:
    }

    \vspace{0.2cm}

    \digitbox Quantidade de provas recolhidas confere com presentes
    \vspace{0.2cm}

    \digitbox Todas as folhas extras foram recolhidas
    \vspace{0.2cm}

    \digitbox Não houve irregularidades
    \vspace{0.2cm}

    \digitbox Houve irregularidades registradas na seção de ocorrências
    \vspace{0.2cm}

    \vspace{2cm}
\end{multicols}

% ====================================================
% ASSINATURAS
% ====================================================

{% if exam.features.closing_city.is_some() %}
\begin{center}
    {{ exam.features.closing_city.as_ref().unwrap() }}, {{ exam.date_extenso() }}.
\end{center}
{% endif %}

{% if exam.features.second_proctor %}
\begin{center}
    \begin{minipage}{0.45\textwidth}
        \centering

        \vspace{0.8cm}

        \rule{6cm}{0.4pt}\\[0.15cm]
        Fiscal Titular\\[0.15cm]
        {\small Nome legível: \dotted{4.5cm}}
    \end{minipage}
    \hfill
    \begin{minipage}{0.45\textwidth}
        \centering

        \vspace{0.8cm}

        \rule{6cm}{0.4pt}\\[0.15cm]
        Fiscal Auxiliar\\[0.15cm]
        {\small Nome legível: \dotted{4.5cm}}
    \end{minipage}
\end{center}
{% else %}
\begin{center}
    \begin{minipage}{0.6\textwidth}
        \centering

        \vspace{0.8cm}

        \rule{7cm}{0.4pt}\\[0.15cm]
        Aplicador\\[0.15cm]
        {\small Nome legível: \dotted{5cm}}
    \end{minipage}
\end{center}
{% endif %}

\end{document}
